\documentclass[10pt,a4paper]{amsart}
\usepackage[margin=19mm]{geometry}
\usepackage{amsmath,amssymb,mathtools}
\usepackage[scr=rsfs,cal=euler]{mathalfa}
\usepackage{xfrac}
\usepackage[unicode,hidelinks,bookmarks=false]{hyperref}
\newcommand{\ie}{i.e.\ }
\newcommand{\cf}{cf.\ }
\newcommand{\resp}{respectively\ }
\newcommand{\Subsection}{\subsection}
\providecommand{\nobreakdash}{\nobreak\hbox{-}}
\setlength{\emergencystretch}{2em}
\multlinegap0pt
\usepackage{amssymb}
\newtheorem{theorem}{Theorem}
\title{Inverse problem for a multi-term time-fractional diffusion equation with the Caputo derivatives}
\author{Ravshan Ashurov, Damir Shamuratov}
\date{Source: arXiv:2603.01833v1 (CC BY 4.0)}
\begin{document}
\maketitle

\noindent\emph{Source and licence.} Inverse problem for a multi-term time-fractional diffusion equation with the Caputo derivatives. Version arXiv:2603.01833v1 (CC BY 4.0), distributed by arXiv under the Creative Commons Attribution 4.0 International licence, http://creativecommons.org/licenses/by/4.0/, as recorded in the arXiv licence metadata for this exact version and shown on the abstract page https://arxiv.org/abs/2603.01833v1. Full licence notice: \texttt{licenses/arxiv-2603.01833-CC-BY-4.0.txt}.\par\smallskip

\noindent\textit{Editorial scope.} Counted unit: the article's theorem theorem (source0.tex lines 1177--1180). The article's complete original proof of it is delivered with it (source0.tex lines 1182--1273, one \texttt{proof} environment of 92 lines). No mathematics is written, repaired or re-derived: the statement and the proof are the article's own lines, in the article's own notation, and no retained line is substituted or re-flowed.  Retained uncounted as the source-local context the proof needs: lines 111--113; lines 121--123; lines 138--140; lines 141--143; lines 815--818.  Nothing was omitted from any retained span, so the package carries no omission record.  No retained line contains a citation, so the wrapper carries no bibliography and no bibliography entry.  No retained line contains a figure environment, an image-inclusion command or drawing code, so no image is delivered and the package declares no asset dependency.  Editorial formatting only, in addition: an AMS \texttt{amsart} preamble that restates the article's own declarations and macros used by the retained lines, namely the environments \texttt{theorem} on separate counters, together with the standard packages those lines need.  The retained environments print in this document as Theorem 1 in order of appearance, whereas the article prints the same items with the numbering of its own full document.  The complete native source of the article as distributed in the arXiv e-print for this exact version is bundled unchanged as \texttt{sources/195579/source0.tex}.
\begin{equation}\label{1.1}
 \sum_{j=1}^M q_j  \partial_t^{\rho_j} u(x,t) + A(x,D)u(x,t) = f(x)g(t), \quad x \in \Omega, \; 0 < t \le T.
\end{equation}

\begin{equation}\label{1.4}
    u(x, t_0) = \Psi(x), \quad x \in \Omega,
\end{equation}

\begin{equation}\label{2.1}
    A(x, D)v(x) = \lambda v(x), \quad x \in \Omega,
\end{equation}

\begin{equation}\label{2.2}
    B_j v(x) = 0, \quad j = 1, \dots, l, \quad x \in \partial \Omega.
\end{equation}

\begin{equation}\label{5.5}
T_k(t) = \varphi_k \left[ 1 - \lambda_k t^{\rho_1} E_{ \rho', \rho_1 + 1} \left( -\lambda_k t^{\rho_1}, * \right) \right]
+ f_kb_{k,\rho_1}(t).
\end{equation}

\begin{theorem}
Let the functions $\varphi(x)$ and $\Psi(x)$ be continuous in the closed domain $\Omega$, and let $K_{0,\rho_1}$ be empty.
Then the problem \eqref{1.1}--\eqref{1.4} admits at most one classical solution $(u(x,t),f(x))$.
\end{theorem}

\begin{proof}
Assume that, under the conditions of the theorem, there exist two pairs of solutions,
$(u_1(x,t),f_1(x))$ and $(u_2(x,t),f_2(x))$.
We prove that
\[
u(x,t)=u_1(x,t)-u_2(x,t)\equiv 0,
\qquad
f(x)=f_1(x)-f_2(x)\equiv 0.
\]

Since the problem under consideration is linear, the pair $(u(x,t),f(x))$ satisfies the following problem:
\begin{equation}\label{6.1}
\sum_{j=1}^M q_j \partial_t^{\rho_j} u(x,t) + A(x,D)u(x,t) = f(x)g(t),
\quad x \in \Omega,\; 0 < t \le T,
\end{equation}
\begin{equation}\label{6.2}
u(x,0) = 0,
\quad x \in \Omega,
\end{equation}
\begin{equation}\label{6.3}
B_j u(x,t) = \sum_{|\alpha|\le m_j} b_{\alpha,j}(x) D^\alpha u(x,t) = 0,
\end{equation}
\[
0 \le m_j \le m-1,\quad j=1,\dots,l,\quad x \in \partial\Omega,\; 0 < t \le T,
\]
\begin{equation}\label{6.4}
u(x,t_0)=0,
\quad x \in \Omega.
\end{equation}

Let $v_k(x)$ be an arbitrary eigenfunction of the spectral problem \eqref{2.1}--\eqref{2.2}
corresponding to the eigenvalue $\lambda_k$.
Consider the function
\begin{equation}\label{6.5}
w_k(t)=\int_{\Omega} u(x,t)v_k(x)\,dx.
\end{equation}

From \eqref{6.1}, we obtain
\[
\sum_{j=1}^M q_j \partial_t^{\rho_j} w_k(t)
=
\sum_{j=1}^M q_j \int_{\Omega} \partial_t^{\rho_j} u(x,t)v_k(x)\,dx
\]
\[
=
- \int_{\Omega} A(x,D)u(x,t)v_k(x)\,dx
+ g(t)\int_{\Omega} f(x)v_k(x)\,dx.
\]

Using the definition of eigenfunctions, we see that
\[
\sum_{j=1}^M q_j \partial_t^{\rho_j} w_k(t)
=
- \int_{\Omega} u(x,t)A(x,D)v_k(x)\,dx + f_k g(t)
=
- \lambda_k w_k(t) + f_k g(t).
\]

Consequently, $w_k(t)$ satisfies the Cauchy-type problem
\[
\sum_{j=1}^M q_j \partial_t^{\rho_j} w_k(t) + \lambda_k w_k(t)
=
f_k g(t),
\qquad
w_k(0)=0,
\]
\begin{equation}\label{6.6}
w_k(t_0)=0.
\end{equation}
By \eqref{5.5}, the solution of this problem has the form
\[
w_k(t)=f_k b_{k,\rho_1}(t).
\]
Using the additional condition \eqref{6.6}, we obtain
\begin{equation}\label{6.7}
w_k(t_0)=f_k b_{k,\rho_1}(t_0)=0.
\end{equation}
Since $K_{0,\rho_1}$ is empty, we have
\[
b_{k,\rho_1}(t_0)\neq 0
\quad \text{for all } k.
\]
Therefore, $f_k=0$ for all $k$.
Hence, $w_k(t)\equiv 0$ for all $k$.
By the completeness of the system of eigenfunctions $v_k$, it follows that
\[
f(x)\equiv 0,
\qquad
u(x,t)\equiv 0.
\]
This proves the uniqueness of the solution to the inverse problem \eqref{1.1}--\eqref{1.4}.
\end{proof}

\end{document}
